\selectlanguage{english}
\begin{abstract}
Municipal governments often need to draft Bills with little time and limited access to comparable legislative evidence. This paper presents a reproducible pipeline to discover Brazilian SAPL instances, extract municipal Bills, translate legislative summaries into operational actions, retrieve similar proposals with semantic search, and relate recurrent policies to official municipal indicators. In a frozen execution dated November 12, 2025, the pipeline found 1,259 SAPL instances, achieved valid extraction in 322 of them, and collected 220,065 Bills. The text-to-action translator was trained on 1,000 synthetic pairs derived from real summaries and obtained 84\% BERTScore on the validation split. We report the method, the resulting dataset, an exploratory case study on violence against women, and the main current limitations of the approach.
\end{abstract}

\noindent\textbf{Keywords}: municipal legislation; SAPL; semantic search; policy recommendation; open data; reproducibility

\selectlanguage{brazil}
\begin{resumo}
Gestores municipais frequentemente precisam redigir Projetos de Lei (PLs) com pouco tempo e pouca visibilidade sobre soluções já propostas em outros municípios. Este artigo apresenta um pipeline reprodutível para descobrir instâncias brasileiras do SAPL, extrair PLs municipais, traduzir ementas em ações operacionais, recuperar propostas semelhantes por busca semântica e relacionar políticas recorrentes a indicadores oficiais municipais. Em uma execução congelada em 12 de novembro de 2025, o pipeline encontrou 1.259 instâncias do SAPL, obteve extração válida em 322 delas e coletou 220.065 PLs. O tradutor \textit{ementa} $\rightarrow$ \textit{ação} foi treinado com 1.000 pares sintéticos derivados de ementas reais e atingiu BERTScore de 84\% no conjunto de validação. O artigo descreve o método, o dataset resultante, um estudo de caso exploratório sobre violência contra a mulher e as principais limitações atuais da abordagem.
\end{resumo}

\noindent\textbf{Palavras-chave}: legislação municipal; SAPL; busca semântica; recomendação de políticas; dados abertos; reprodutibilidade
